\section{Thiết kế demo thực nghiệm}
\label{sec:demo_design}

\subsection{Mục tiêu và triết lý thiết kế}

Demo có tên \textbf{Mini Benchmark Evaluator}. Điểm khác biệt cốt lõi so với một
bài ``so sánh model'' thông thường: \emph{nhóm không xếp hạng model như mục
tiêu cuối}. Thay vào đó, nhóm dùng nhiều model như những \emph{đầu dò} để
\textbf{soi chính phép đo} --- benchmark có đo đúng thứ nó tuyên bố, điểm số có ý
nghĩa, và thứ hạng có ổn định không. Đây là hiện thực hoá trực tiếp tinh thần
\emph{Measuring what Matters}~\cite{measuring_what_matters} và bộ tiêu chí
\emph{BetterBench}~\cite{betterbench}.

Một quyết định thiết kế then chốt: thay vì \emph{tự chế} một benchmark đồ chơi
(dễ nhưng không thuyết phục), nhóm lấy \textbf{hai benchmark công khai có
thật} --- MMLU và GSM8K --- rồi áp bộ công cụ phân tích phê phán của tutorial lên
chúng. Nhờ vậy mọi hiện tượng quan sát được (bão hoà, nhiễm dữ liệu, độ giòn của
metric) là \emph{thật}, không phải dàn dựng.

\subsection{Kiến trúc hệ thống}

Hệ thống tách bạch theo trách nhiệm để dễ tái lập và kiểm tra (Hình~\ref{fig:arch}
mô tả luồng):

\begin{description}
  \item[\texttt{config.py}] Khai báo tập con dữ liệu (kích thước, seed) và danh
    sách \emph{systems} (6 model + 2 baseline). Thêm/bớt model = sửa một danh sách.
  \item[\texttt{data.py}] Tải MMLU + GSM8K từ Hugging Face, lấy tập con cố định
    theo seed, ghi ra \texttt{data/*.jsonl}.
  \item[\texttt{perturb.py}] Sinh ``bản sinh đôi'' nhiễu giữ nhãn của câu MMLU
    (đảo vị trí đáp án; chèn distractor ``None of the above'') để đo robustness.
  \item[\texttt{models.py}] \texttt{HFRunner} chạy mô hình Hugging Face ở chế độ
    4-bit; \texttt{BaselineRunner} (random/majority); \texttt{FakeRunner} (test
    pipeline không GPU). Mọi lời gọi đi qua một \emph{cache} theo
    \texttt{sha256(model\_id, prompt)}.
  \item[\texttt{evaluate.py}] Vòng lặp đánh giá: nạp \emph{từng} model một, chạy
    hết item, giải phóng GPU, sang model kế tiếp; ghi \texttt{results/per\_item.csv}.
  \item[\texttt{analysis.py}] Từ \texttt{per\_item.csv} sinh toàn bộ bảng số và
    biểu đồ.
\end{description}

\begin{figure}[t]
  \centering
  % Thay bằng sơ đồ kiến trúc nếu có; tạm mô tả bằng pipeline văn bản.
  \fbox{\parbox{0.92\linewidth}{\centering\small
    \texttt{config} $\to$ \texttt{data}/\texttt{perturb} $\to$
    \texttt{evaluate} (nạp 1 model $\to$ cache/generate $\to$ unload) $\to$
    \texttt{per\_item.csv} $\to$ \texttt{analysis} $\to$ bảng + hình}}
  \caption{Luồng xử lý của Mini Benchmark Evaluator. Mỗi model được nạp và giải
  phóng độc lập nên cả 6 model chạy gọn trên một GPU T4 free.}
  \label{fig:arch}
\end{figure}

\subsection{Liên hệ với sáu paper}

Thiết kế demo được neo có chủ đích vào cả sáu paper: phủ \emph{hẹp} của benchmark
(Raji), tách \emph{metric} khỏi \emph{construct} và dùng baseline (Measuring what
Matters), đủ README/metric/baseline/visualize/tái lập (BetterBench), gộp nhiều
nhóm task thành một suite nhỏ (GLUE), thêm perturbation đo robustness khi đề bão
hoà (SuperGLUE), và chấm bằng \emph{trích xuất + so giá trị} thay vì so chuỗi thô
(tinh thần execution-based của SWE-bench).
